% chapters/01/section-04.tex
\section{Konteks Kurikulum OBE dan Sistem Asesmen}
\label{sec:kurikulum-obe}

Bagian ini menguraikan bagaimana capaian yang telah Anda baca sebelumnya
diterjemahkan menjadi penilaian. Memahami bagian ini sejak awal akan membuat
Anda bekerja dengan arah yang jelas: Anda mengetahui bukti apa yang harus
dihasilkan, kapan harus dihasilkan, dan dengan ukuran apa dinilai.

\subsection{Ranah Taksonomi Bloom dan Tuntutan Tiap Tingkat}

Setiap Sub-CPMK pada mata kuliah ini diberi tingkat ranah kognitif menurut
taksonomi Bloom, mulai dari C2 (memahami) sampai C5 (mengevaluasi). Tingkat
tersebut bukan hiasan, melainkan penanda seberapa jauh Anda dituntut menguasai
materi. Tabel~\ref{tab:bloom} menjelaskan tuntutan setiap tingkat beserta
contohnya pada mata kuliah ini.

\begin{table}[htbp]
\centering
\caption{Tingkat ranah Bloom, tuntutannya, dan contoh pada mata kuliah ini}
\label{tab:bloom}
{\setlength{\tabcolsep}{4pt}\footnotesize%
\begin{tabularx}{\textwidth}{@{}C{1.0cm} L{3.4cm} L{4.3cm} Y@{}}
\toprule
\headrow \textbf{Tkt} & \textbf{Kata kerja khas} & \textbf{Tuntutan} & \textbf{Contoh pada mata kuliah ini} \\
\midrule
C2 & menjelaskan, mengidentifikasi, membedakan &
Menunjukkan pemahaman dengan kata-kata sendiri, tanpa menuntut karya baru. &
Menjelaskan urutan tahapan \istilah{game loop} dan peran pembatasan laju
bingkai (\SUBCPMK{1.1}). \\
\addlinespace
C3 & menulis, mengimplementasikan, mengonfigurasi &
Menghasilkan karya yang berjalan sesuai maksud, dengan mengikuti konsep yang
sudah dipelajari. &
Menulis kode gerak berbasis \istilah{delta time} pada Pygame
(\SUBCPMK{1.2}); mengonfigurasi \istilah{Collider} pada Unity
(\SUBCPMK{3.2}). \\
\addlinespace
C4 & menganalisis, memilih, menyesuaikan, memperbaiki &
Menemukan penyebab kegagalan atau memilih di antara beberapa cara, dengan
alasan yang dapat dipertanggungjawabkan. &
Memilih metode deteksi tabrakan yang sesuai dan memperbaiki tabrakan yang
lolos deteksi (\SUBCPMK{1.4}); menganalisis pengaruh heuristik pada A*
(\SUBCPMK{2.3}). \\
\addlinespace
C5 & mengevaluasi, menilai, merumuskan perbaikan &
Menilai mutu hasil kerja sendiri berdasarkan kriteria, lalu merumuskan
perbaikan yang ditindaklanjuti. &
Menilai perilaku NPC hasil gabungan mesin keadaan dan A* terhadap skenario
permainan (\SUBCPMK{2.4}); mengevaluasi prototipe melalui
\istilah{playtesting} (\SUBCPMK{4.3}). \\
\bottomrule
\end{tabularx}}
\end{table}

Tuntutan meningkat dari blok ke blok. Blok I dan Blok II
banyak menuntut C3 dan C4 (menghasilkan dan menganalisis), sedangkan Blok II
dan Blok IV menutup dengan C5 (mengevaluasi). Pola ini menjelaskan mengapa
laporan \istilah{playtest} dan refleksi menjadi bagian penilaian, bukan
pelengkap.

\subsection{Asesmen Autentik dan Portofolio}

Asesmen pada mata kuliah ini bersifat autentik: yang dinilai adalah artefak
nyata hasil pekerjaan Anda, bukan kemampuan mengingat. Bentuk artefaknya
berupa kode program yang dapat dijalankan, level yang dapat dimainkan,
laporan pengujian, dan lembar refleksi.

Praktikum dikerjakan setiap minggu secara individual dan dikumpulkan melalui
repositori Git pribadi. Setiap tugas memuat kode program, tangkapan layar
hasil, serta catatan singkat mengenai kendala yang dihadapi dan cara
penyelesaiannya. Dosen menilai enam sampel portofolio secara acak sebagai
dasar nilai portofolio. Cara ini dipilih dengan sengaja: karena sampelnya
acak, mengerjakan seluruh rangkaian tugas secara berkelanjutan menjadi jauh
lebih menguntungkan daripada mengerjakan sebagian tugas dengan sangat baik.

Umpan balik diberikan setiap minggu, bukan hanya pada akhir blok. Manfaatkan
umpan balik tersebut selagi materinya masih dekat, karena memperbaiki kode
yang Anda tulis minggu ini jauh lebih mudah daripada memperbaiki kode yang
sama tiga minggu kemudian.

\subsection{Komponen Asesmen dan Waktu Pelaksanaannya}

Penilaian pada mata kuliah ini berasal dari enam komponen yang tersebar
sepanjang semester.

\begin{itemize}
  \item \textbf{Tugas praktikum dan portofolio} (dikerjakan sepanjang
        semester). Dua belas tugas praktikum pada minggu 1--8 dan 11--14
        dikumpulkan melalui repositori Git, masing-masing berisi kode,
        tangkapan layar, dan catatan kendala. Enam di antaranya dinilai sebagai
        sampel. Komponen ini mengukur \CPMK{1} sampai \CPMK{4}.
  \item \textbf{Kuis 1} (minggu 5) atas Materi 1--4 dan \textbf{Kuis 2}
        (minggu 8) atas Materi 6--7. Keduanya berupa tes singkat 15--20 menit
        yang mengukur \CPMK{1} dan \CPMK{2} secara formatif.
  \item \textbf{Mini proyek AI} (dikumpulkan dan dipresentasikan pada minggu
        9). Dinilai dari kode NPC bermesin keadaan dan pencarian jalur A*,
        laporan \istilah{playtest}, dan presentasi kelompok. Komponen ini
        mengukur \CPMK{2}.
  \item \textbf{Ujian Tengah Semester} (minggu 10) atas \CPMK{1} dan
        \CPMK{2}, berupa tes tulis konsep dan uji kode perbaikan kesalahan pada
        program Pygame dalam batas waktu.
  \item \textbf{Proyek akhir} (dikumpulkan pada minggu 15--16). Dinilai dari
        berkas terpasang (\istilah{build}), kode sumber, laporan pengembangan,
        presentasi, dan uji fungsional oleh dosen. Komponen ini mengukur
        \CPMK{3} dan \CPMK{4}.
  \item \textbf{Partisipasi dan kehadiran} (sepanjang semester), dinilai
        melalui observasi kehadiran dan keaktifan pada praktikum, diskusi,
        serta kegiatan \istilah{playtest}.
\end{itemize}

Bobot setiap komponen, matriks bobot terhadap CPMK, rubrik penilaian, dan
ketentuan konversi nilainya diuraikan pada Lampiran~A.

\subsection{Kriteria Ketuntasan dan Ambang Ketercapaian}

Anda dinyatakan lulus mata kuliah ini apabila memenuhi tiga ketentuan
sekaligus:
\begin{enumerate}
  \item memperoleh nilai akhir minimal \textbf{60} (huruf C);
  \item memiliki tingkat kehadiran minimal \textbf{75\%} dari seluruh
        pertemuan;
  \item menyerahkan seluruh keluaran wajib berupa mini proyek AI dan proyek
        akhir.
\end{enumerate}
Nilai akhir dikonversi menjadi huruf dengan rentang A untuk nilai 80 ke atas,
B untuk 70--79, C untuk 60--69, D untuk 50--59, dan E untuk nilai di bawah 50.

Selain kriteria untuk mahasiswa, terdapat pula ambang ketercapaian yang
dipakai dosen untuk menilai keberhasilan penyelenggaraan mata kuliah, yaitu
sekurang-kurangnya 75\% mahasiswa memperoleh nilai 70 atau lebih pada setiap
CPMK. Ambang ini adalah alat pemantauan mutu, bukan syarat tambahan bagi
nilai Anda. Apabila suatu CPMK belum mencapai ambang tersebut, dosen mencatat
penyebabnya---apakah pada materi, metode, atau instrumen asesmennya---lalu
menyesuaikan pelaksanaan pada semester berikutnya.

\subsection{Beban Belajar Tiga Satuan Kredit Semester}

Bobot \BobotSKS\ bukan sekadar penanda jumlah pertemuan. Mengacu pada
ketentuan beban belajar, setiap satu satuan kredit semester setara dengan 50
menit tatap muka, 60 menit belajar terstruktur, dan 60 menit belajar mandiri
setiap minggu. Rinciannya diperlihatkan pada Tabel~\ref{tab:alokasi-buku}.

\begin{table}[htbp]
\centering
\caption{Bentuk kegiatan dan alokasi waktu selama enam belas minggu}
\label{tab:alokasi-buku}
{\setlength{\tabcolsep}{4pt}\footnotesize%
\begin{tabularx}{\textwidth}{@{}L{3.2cm} Y C{2.5cm} C{3.0cm}@{}}
\toprule
\headrow \textbf{Bentuk} & \textbf{Kegiatan} & \textbf{Per minggu} & \textbf{Total 16 minggu} \\
\midrule
Tatap muka (TM) &
Kuliah interaktif, demonstrasi kode, dan praktikum terbimbing di
laboratorium. &
150 menit & 2.400 menit (40 jam) \\
\addlinespace
Belajar terstruktur (BT) &
Tugas praktikum, latihan terarah, tugas membaca, dan pengerjaan proyek
terjadwal. &
180 menit & 2.880 menit (48 jam) \\
\addlinespace
Belajar mandiri (BM) &
Menelusuri dokumentasi resmi, bereksperimen dengan kode, dan menyelesaikan
proyek. &
180 menit & 2.880 menit (48 jam) \\
\midrule
\textbf{Total} & & \textbf{510 menit} & \textbf{8.160 menit (136 jam)} \\
\bottomrule
\end{tabularx}}
\tabsource{Sumber: RPS mata kuliah \NamaMK.}
\end{table}

Angka tersebut menjelaskan satu hal yang sering diremehkan: waktu tatap muka
hanya menempati sekitar tiga puluh persen dari seluruh beban belajar. Sisanya
berada di tangan Anda, dalam bentuk mengerjakan praktikum dan mempelajari
dokumentasi. Mata kuliah pemrograman tidak dapat dikuasai hanya dengan hadir
dan menyimak; keterampilan menulis program terbentuk melalui latihan yang
cukup banyak dan berulang.
